\chapter[The microscopic action]{Matrix Coordinates and the Microscopic Action}
\label{ch:matrix-action}

\paragraph{Learning outcomes.}
After this chapter the reader should be able to:
\begin{enumerate}
  \item write the non-commutative coordinate algebra
        \([\hat X^\mu,\hat X^\nu]=i\Theta^{\mu\nu}\) and state how
        \(\Theta^{\mu\nu}\) supplies an ultraviolet cutoff;
  \item write both microscopic actions used in the NAMCG notes---the
        pregeometric matrix action and the covariant condensate-field
        action---and say what each is for;
  \item explain why the order parameter must carry an internal index
        \(I\) and a tangent index \(A\);
  \item expand the Mexican-hat matrix potential about its vacuum and
        identify the radial and Goldstone directions;
  \item enter Chapter~4 knowing which object will condense into a
        vierbein.
\end{enumerate}

Chapter~2 supplied geometry as language. This chapter supplies the
field. After this point the metric is not an independent ingredient.
It is something the field must produce.

\section{The scalar failure, restated as a field-content problem}
\label{sec:ch3-scalar-failure}

A single complex scalar \(\phi\) transforming as a singlet under every
tangent-space representation can condense. It can eat one Goldstone
phase into a gauge field. It cannot furnish a local Lorentz frame.
The infrared excitations of such a condensate are, at best,
\begin{itemize}
  \item a radial Higgs scalar,
  \item a massive or massless vector, if a gauge field is present,
  \item phonons that admit an acoustic metric rewrite.
\end{itemize}
There is no symmetric trace-free tensor \(T_{\mu\nu}\) hiding in
\(\phi\). Seeking that tensor in a more elaborate potential for the
same \(\phi\) is a representation-theoretic error, not a calculational
one.

NAMCG therefore promotes the order parameter to a rectangular (or
square, after a choice of basis) matrix field
\begin{equation}
  \label{eq:Phi-field}
  \Phi^I{}_A(x),
\end{equation}
where \(I=1,\ldots,N\) labels an internal gauge representation and
\(A=0,1,2,3\) labels an emergent tangent space. The field is a
linear map from the tangent space of an auxiliary frame into the
internal space, or the reverse, according to index placement. The
placement used in this book is: \(I\) up, \(A\) down on \(\Phi\),
unless a formula from the source notes is being quoted verbatim.

The rest of the chapter writes an action for \(\Phi\) and for a
closely related pregeometric object, the matrix coordinates
\(\hat X^\mu\).

\section{Two presentations, and why both appear}
\label{sec:ch3-two-presentations}

The original notes use two microscopic languages. They are not
identical, and treating them as interchangeable by slogan would
produce a confused book.

\paragraph{Presentation A: pregeometric matrix coordinates.}
There is, at the outset, no manifold. The dynamical objects are
Hermitian operators \(\hat X^\mu\) acting on a Hilbert space
\(\mathcal{H}\). Spacetime points, if they appear at all, appear as
eigenvalue data of a complete commuting subset of these operators, or
as a derived continuum after condensation.

\paragraph{Presentation B: condensate field on an auxiliary background.}
There is an auxiliary manifold with a non-dynamical metric
\(g^{(0)}_{\mu\nu}\). The dynamical object is the field
\(\Phi^I{}_A(x)\). The physical metric \(g_{\mu\nu}\) is not
\(g^{(0)}_{\mu\nu}\). The auxiliary metric is a scaffold for writing
derivatives and heat kernels. Chapter~8 uses this scaffold.

This book treats Presentation~A as the ontological starting point and
Presentation~B as the background-field language appropriate after a
condensate has selected a notion of local coordinates. Chapter~4 is
the bridge. Until that bridge is built, the two actions below are
written separately and not added together.

\begin{assumption}[Auxiliary background is non-dynamical]
  \label{ass:aux-metric}
  Whenever Presentation~B is in use, \(g^{(0)}_{\mu\nu}\) is a fixed
  scaffold, usually taken flat,
  \(g^{(0)}_{\mu\nu}=\eta_{\mu\nu}\). It is not varied in the
  Einstein--Hilbert sense. Varying the \emph{emergent} metric
  \(g_{\mu\nu}\) defined in Chapter~4 is a different operation.
\end{assumption}

\section{Presentation A: matrix coordinates and ultraviolet cutoff}
\label{sec:ch3-matrices}

\begin{definition}[Matrix coordinates]
  \label{def:matrix-coords}
  Let \(\mathcal{H}\) be a separable Hilbert space and let
  \(\hat X^\mu\), \(\mu=0,1,2,3\), be Hermitian operators on
  \(\mathcal{H}\). The fundamental kinematic relation of Presentation~A
  is
  \begin{equation}
    \label{eq:NC}
    [\hat X^\mu,\hat X^\nu]=i\Theta^{\mu\nu},
  \end{equation}
  where \(\Theta^{\mu\nu}=-\Theta^{\nu\mu}\) is an antisymmetric
  tensor with dimensions of length squared. If \(\Theta^{\mu\nu}\) is
  a multiple of the identity on \(\mathcal{H}\), it is a
  \(c\)-number background. In the fully dynamical theory it may
  itself be an operator; the notes treat the leading infrared
  analysis with a \(c\)-number \(\Theta^{\mu\nu}\).
\end{definition}

The scale
\begin{equation}
  \label{eq:LUV-def}
  \LUV\sim\Theta^{-1/2},
  \qquad
  \Theta
  \equiv
  \sqrt{\tfrac12\Theta_{\mu\nu}\Theta^{\mu\nu}},
\end{equation}
is the ultraviolet cutoff of the theory. Two operators that fail to
commute cannot be simultaneously diagonalized, so the notion of a
sharply localized event \((x^0,x^1,x^2,x^3)\) is meaningless below
\(\LUV^{-1}\). This is the NAMCG replacement for a hard lattice.

\begin{remark}[What \(\Theta^{\mu\nu}\) is not]
  Equation~\eqref{eq:NC} is not yet a solution of the cosmological-
  constant problem, not yet a black-hole core, and not yet a graviton.
  It is a kinematic regulator. Claims that begin ``because spacetime
  is non-commutative\ldots'' and end at galactic rotation curves are
  skipping the work of Chapters~4--14.
\end{remark}

A Yang--Mills-type matrix action for these coordinates, including a
cubic structure-constant coupling to a background condensate tensor
\(C_{\mu\nu\lambda}\), is
\begin{equation}
  \label{eq:Smatrix}
  \begin{aligned}
    S_{\mathrm{matrix}}
    &=-\frac{1}{4g_{\mathrm{YM}}^2}
      \mathrm{Tr}\bigl(
        [\hat X^\mu,\hat X^\nu][\hat X_\mu,\hat X_\nu]
      \bigr)
      +\frac{1}{3!}
      \mathrm{Tr}\bigl(
        C_{\mu\nu\lambda}[\hat X^\mu,\hat X^\nu]\hat X^\lambda
      \bigr)\\
    &\qquad
      +S_{\mathrm{matter}}(\Psi,\hat X^\mu).
  \end{aligned}
\end{equation}
The first term is the IKKT / reduced-Yang--Mills kinetic term. On a
background satisfying~\eqref{eq:NC} it is proportional to
\(\mathrm{Tr}(\Theta_{\mu\nu}\Theta^{\mu\nu})\) plus fluctuation
pieces. The cubic term selects a preferred vacuum condensation
structure; its detailed index symmetries are constrained by the
requirement that the vacuum preserve a Lorentzian interpretation
after Chapter~4.

Matter fields \(\Psi\) in \(S_{\mathrm{matter}}\) are functions of the
matrices \(\hat X^\mu\) in the sense of the Weyl map associated with
\(\Theta^{\mu\nu}\). This book does not need the full star-product
calculus until phenomenology with explicit localized sources appears
in Chapter~14. The existence of such a map is enough for now.

\section{Presentation B: the covariant condensate action}
\label{sec:ch3-Phi-action}

Once local coordinates \(x^\mu\) have been introduced as labels on an
auxiliary manifold, the dynamical variable is the matrix field
\(\Phi^I{}_A(x)\). Let \(A_\mu\) be a connection on the internal
bundle, with covariant derivative
\begin{equation}
  \label{eq:Dmu-Phi}
  D_\mu\Phi=\partial_\mu\Phi-iA_\mu\Phi
\end{equation}
acting on the \(I\) index. The microscopic action used for the
heat-kernel analysis is
\begin{equation}
  \label{eq:Smicro}
  S_{\mathrm{micro}}[\Phi,A]
  =\int\mathrm{d}^4x\,\sqrt{-g^{(0)}}\,
   \mathrm{Tr}\Bigl(
    -\,g^{(0)\mu\nu}(D_\mu\Phi)^\dagger D_\nu\Phi
     -V(\Phi^\dagger\Phi)
   \Bigr).
\end{equation}
The trace is over internal indices. Tangent indices on \(\Phi\) are
contracted through the potential and through the identification of
\(A\) with a frame index after condensation; they are not traced
in the same sum.
In the signature \((-,+,+,+)\) of Chapter~2 the kinetic term is written with
the minus sign so that time derivatives have positive energy. This is a
correction to the Presentation-B action as stated; that the matrix action
of Presentation~A reproduces this sign is part of the unproved dictionary
(Assumption~\ref{ass:dictionary}) and remains to be derived.

The potential that triggers condensation is the matrix Mexican hat
\begin{equation}
  \label{eq:V}
  V(\Phi^\dagger\Phi)
  =\frac{\lambda}{4}
   \mathrm{Tr}\bigl(\Phi^\dagger\Phi-v^2\mathbb{I}\bigr)^2.
\end{equation}
Its minima lie on the manifold
\begin{equation}
  \label{eq:vacuum-manifold}
  \Phi^\dagger\Phi=v^2\mathbb{I}
\end{equation}
in the internal space, up to unitary rotations. A representative
vacuum used throughout the book is
\begin{equation}
  \label{eq:VEV}
  \langle\Phi^I{}_A\rangle=v\,\delta^I{}_A,
\end{equation}
which is available when the internal dimension is at least four and a
basis has been chosen that identifies four internal directions with
the tangent directions. If \(N>4\), the remaining internal directions
must either decouple or be locked by additional Yukawa-like couplings.
The source notes assume that a vacuum of the form~\eqref{eq:VEV}
exists. This book records that as an assumption.

\begin{assumption}[Existence of a frame-locking vacuum]
  \label{ass:frame-lock}
  There exists a stable vacuum~\eqref{eq:VEV} (or a locally equivalent
  vacuum) in which four internal directions are identified with the
  four tangent directions. Goldstone modes associated with the broken
  generators are available to be eaten by a connection, as in
  Chapter~10.
\end{assumption}

\section{Index structure and the intended infrared spectrum}
\label{sec:ch3-indices}

Under an internal gauge transformation \(\Omega(x)\) and a local
Lorentz transformation \(\Lambda(x)\), the field is required to
transform as
\begin{equation}
  \label{eq:Phi-transform}
  \Phi^I{}_A(x)
  \;\longmapsto\;
  \Omega^I{}_J(x)\,\Phi^J{}_B(x)\,\Lambda^B{}_A(x).
\end{equation}
Equation~\eqref{eq:Phi-transform} is the transformation law that
makes \(\Phi\) capable of locking the two groups to a diagonal
subgroup. It is a choice of field content. It is not derived from
Presentation~A alone; it is the additional structure that Presentation~B
puts in by hand so that a vierbein can exist.

Fluctuations about~\eqref{eq:VEV} will be decomposed in Chapter~5
into
\begin{equation}
  \label{eq:split-preview}
  \delta\Phi^I{}_A
  =\frac14\eta_{\mu\nu}\delta g^{\mu\nu}\,I
   +A_\mu\gamma^\mu
   +T_{\mu\nu}.
\end{equation}
The only fact needed here is that the last piece exists because
\(\Phi\) has two indices that can be symmetrized and made traceless
after the vacuum has identified \(I\) with \(A\). A scalar has no
such piece.

\section{Equations of motion from Presentation B}
\label{sec:ch3-eom}

Varying~\eqref{eq:Smicro} with respect to \(\Phi^\dagger\) gives
\begin{equation}
  \label{eq:Phi-eom}
  D_\mu\bigl(\sqrt{-g^{(0)}}\,g^{(0)\mu\nu}D_\nu\Phi\bigr)
  -\sqrt{-g^{(0)}}\,
   \frac{\lambda}{2}\Phi
   \bigl(\Phi^\dagger\Phi-v^2\mathbb{I}\bigr)
  =0.
\end{equation}
Varying with respect to \(A_\mu\) produces a matrix current
conservation law
\begin{equation}
  \label{eq:gauge-eom}
  D_\mu
  \bigl(
    \Phi(D^\mu\Phi)^\dagger-(D^\mu\Phi)\Phi^\dagger
  \bigr)
  =0,
\end{equation}
up to the conventional factor of \(i\) in the current definition.
These equations determine the condensate profile on the auxiliary
background. They do not yet determine the Einstein tensor of the
emergent metric. That requires integrating out fluctuations, which is
Chapter~8, after the emergent metric has been defined in Chapter~4
and the tensor mode identified in Chapter~5.

\section{Relating the two presentations}
\label{sec:ch3-relate}

A complete derivation from~\eqref{eq:Smatrix} to~\eqref{eq:Smicro}
would construct, from the spectra of \(\hat X^\mu\), a commutative
manifold, a Weyl map, and a slowly varying matrix-valued field on that
manifold whose covariant derivative reproduces the commutator
\([\hat X^\mu,\hat\Phi]\). That derivation is not supplied in full in
the source notes. The working relation used in this book is the
following.

\begin{assumption}[Infrared dictionary]
  \label{ass:dictionary}
  In a condensate phase in which the matrices \(\hat X^\mu\) are,
  to leading order, commutative coordinates \(x^\mu\mathbb{I}\) plus
  fluctuations, there exists a slowly varying field
  \(\Phi^I{}_A(x)\) such that
  \begin{equation}
    \label{eq:dictionary}
    [\hat X_\mu,\Phi]
    \;\longleftrightarrow\;
    i\Theta_{\mu\nu}D^\nu\Phi
    \qquad\text{plus higher-derivative corrections in \(\Theta\).}
  \end{equation}
  Under this dictionary the quadratic commutator term
  in~\eqref{eq:Smatrix} is required to generate the kinetic term
  in~\eqref{eq:Smicro}, including its sign in the signature
  \((-,+,+,+)\); this is a matching condition, not an established result.
  \(\LUV\) remains the cutoff of the derivative expansion.
\end{assumption}

Assumption~\ref{ass:dictionary} is the most important unproved bridge
in the microscopic setup. Chapter~17 lists a sharper derivation of
this dictionary as an open problem. Every later heat-kernel formula
is a calculation \emph{within} Presentation~B, contingent on the
dictionary.

\section{Energy scales and coupling constants}
\label{sec:ch3-scales}

The dimensionful quantities introduced so far are
\begin{equation}
  \label{eq:scales}
  \LUV,\qquad v,\qquad \lambda,\qquad g_{\mathrm{YM}},
\end{equation}
together with the auxiliary-background curvature scale, which is set
to zero when \(g^{(0)}=\eta\). Their intended roles are:

\begin{center}
\begin{tabular}{@{}ll@{}}
\toprule
Quantity & Role \\
\midrule
\(\LUV\) & kinematic cutoff from \(\Theta^{\mu\nu}\) \\
\(v\) & condensate scale; appears in the working conditional \(\Geff\) parameterization \\
\(\lambda\) & controls the radial Higgs mass \(m_\sigma\sim\sqrt{\lambda}\,v\) \\
\(g_{\mathrm{YM}}\) & internal gauge coupling; enters the connection
                     that eats Goldstones \\
\bottomrule
\end{tabular}
\end{center}

Chapter~8 gives a conditional truncated heat-kernel identity. Define
\(C_R^L\) as the Lorentzian curvature coefficient only after the
specified sector sum and continuation:
\begin{equation}
  \label{eq:Geff-again}
  C_R^L[\{\mathcal O_s,w_s,{\LUV}_s,\mu_s\},\mathcal S]
  \quad\text{is not evaluated here.}
\end{equation}
The scale hierarchy needed to integrate out the radial mode is
separate from this coefficient matching:
\begin{equation}
  \label{eq:scale-ineq}
  m_\sigma\ll\LUV,
  \qquad
  v\lesssim\LUV.
\end{equation}
These are working hierarchy assumptions, not consequences of a
computed Newton coupling. No Planck-cutoff inference, cancellation,
or numerical value of \(G_{\mathrm{eff}}\) follows before Problem~F
is completed.

\section{What has been assumed, collected}
\label{sec:ch3-assumptions}

The microscopic theory, as used from here on, rests on:

\begin{enumerate}
  \item the coordinate algebra~\eqref{eq:NC} with cutoff~\eqref{eq:LUV-def};
  \item the matrix action~\eqref{eq:Smatrix} at the pregeometric level;
  \item the condensate action~\eqref{eq:Smicro}--\eqref{eq:V} after
        coordinates exist;
  \item Assumption~\ref{ass:aux-metric}, that \(g^{(0)}\) is a scaffold;
  \item Assumption~\ref{ass:frame-lock}, that a vacuum
        \(\langle\Phi^I{}_A\rangle=v\delta^I{}_A\) exists;
  \item Assumption~\ref{ass:dictionary}, the infrared commutator
        dictionary.
\end{enumerate}

Everything in Chapters~4--11 is a calculation from this list, or an
additional assumption announced as such. The list is long enough that
the theory is not ``derived from nothing.'' It is derived from a
definite microscopic ansatz. That is still a stronger starting point
than inserting \(g_{\mu\nu}\) by hand, provided the later chapters
do the work advertised in Chapter~1.

\section{Exercises}
\label{sec:ch3-exercises}

\begin{enumerate}
  \item \textbf{Dimensions.}
    Assign mass dimensions to \(\hat X^\mu\), \(\Theta^{\mu\nu}\),
    \(\Phi\), \(v\), \(\lambda\), and \(g_{\mathrm{YM}}\) in four
    dimensions, with \(\hbar=c=1\). Confirm~\eqref{eq:LUV-def}.
  \item \textbf{Vacuum of the potential.}
    Show that~\eqref{eq:V} is bounded below and that its minima
    satisfy~\eqref{eq:vacuum-manifold}. Linearize
    \(\Phi=v\,\mathbb{I}+\sigma+i\pi\) in the internal space
    (ignoring tangent structure) and read off \(m_\sigma^2\) and
    \(m_\pi^2\).
  \item \textbf{Why \(N\ge 4\).}
    If the internal dimension is \(N=2\), can a vacuum of the
    form~\eqref{eq:VEV} exist? What infrared representation content
    is then missing?
  \item \textbf{Commutator to derivative.}
    In the Moyal plane with constant \(\Theta^{12}=-\Theta^{21}=\theta\)
    and all other components zero, verify the leading term
    of~\eqref{eq:dictionary} on a plane-wave matrix
    \(\Phi=e^{ik\cdot\hat X}\). Where do higher powers of \(\theta\)
    appear?
  \item \textbf{IKKT vacuum energy.}
    Evaluate the first term of~\eqref{eq:Smatrix} on the
    \(c\)-number background~\eqref{eq:NC}. Show that it is a
    constant proportional to \(\mathrm{Tr}(\Theta_{\mu\nu}\Theta^{\mu\nu})\).
    This constant is a preview of the \(c_0\LUV^4\) problem.
  \item \textbf{Transformation law.}
    Starting from~\eqref{eq:Phi-transform} and the representative
    vacuum~\eqref{eq:VEV}, compute the infinitesimal variation of
    \(\Phi\) under a simultaneous internal rotation and local Lorentz
    transformation that preserve the vacuum. Which generators are
    broken? (This is the algebra of Chapter~10, prepared here.)
  \item \textbf{No graviton in the scalar truncation.}
    Set \(N=1\) and drop the tangent index, reducing \(\Phi\) to a
    complex scalar. Show that the fluctuation spectrum of
    \eqref{eq:Smicro} contains no spin-2 state. This is the
    Chapter~1 ceiling written as a homework problem.
\end{enumerate}

\section{Bridge to Chapter~4}
\label{sec:ch3-bridge}

The field exists, the potential prefers a vacuum
\(\langle\Phi^I{}_A\rangle=v\delta^I{}_A\), and the coordinate algebra
supplies a cutoff. Nothing yet is a metric. Chapter~4 defines the
emergent vierbein as the condensate gradient
\begin{equation}
  \label{eq:vierbein-next}
  e^A_\mu(x)\equiv\langle\partial_\mu\Phi^{IA}(x)\rangle,
\end{equation}
reconstructs \(g_{\mu\nu}=\eta_{AB}e^A_\mu e^B_\nu\), and writes the
Seiberg--Witten map that relates the non-commutative background
\(\Theta^{\mu\nu}\) to the leading correction
\(\mathcal{G}_{\mu\nu}(\langle\hat X\rangle)\). After that chapter,
geometry is an output of the action written here.